% ============================================================
% A Spherical N-Series Area Law for Regular Refinement Families
% Version 2.4 DRAFT FOR AUTHOR REVIEW -- November 10, 2026
%
% Editorial provenance notes:
%  * This draft leads with the manuscript's own results (cubic obstruction,
%    mixed-exponent operators, cap family) and places the 1991--1994 literature on
%    flat-triangle approximation of surface integrals (Georg; Georg--Tausch;
%    Verlinden--Cools; Lyness), which v2.3 did not cite, as background.
%  * Section 6.2 (joint-expansion theorem, cell-coefficient lemma, certified
%    coefficient) folds in the November 10 proof supplement; scripts
%    q2_lemma_symbolic_audit.py, c2prime_enclosure.py, c2prime_crosscheck_numpy.py.
%  * Section 6.4 and Tables 3--4 are new: computed November 10, 2026 with
%    scripts/cubic_amplitude_test_mpfr.py (160-bit MPFR, deterministic, with
%    orientation check); scripts/cubic_amplitude_test.py is the long-double version.
%  * Numerical content of all tables carried over from v2.3 is unchanged.
%  * The comparison table with Georg--Tausch is new; its "this archive"
%    column was computed on November 10, 2026 from the archived routine
%    chordal_triangle in scripts/spherical_nseries_area_suite.py (50 digits).
%  * Exact weights, cubic residual factors and moments were re-derived in
%    exact rational arithmetic (scripts/check_v2_4_draft.py).
%  * Equation and table numbers no longer match v2.1--v2.3.
% ============================================================
\documentclass[11pt]{article}
\usepackage[utf8]{inputenc}
\usepackage[T1]{fontenc}
\usepackage{amsmath,amssymb,amsthm}
\usepackage{booktabs}
\usepackage{xcolor}
\usepackage[margin=1in]{geometry}
\usepackage{hyperref}

\newtheorem{theorem}{Theorem}
\newtheorem{lemma}{Lemma}
\newtheorem{corollary}{Corollary}
\newtheorem{proposition}{Proposition}
\theoremstyle{remark}
\newtheorem{remark}{Remark}

\newcommand{\Ahat}{\widehat{A}}
\newcommand{\Aex}{A_{\mathrm{exact}}}
\newcommand{\R}{\mathbb{R}}

\title{A Spherical N-Series Area Law for Regular Refinement Families\\[2pt]
\large A cubic obstruction under scale-dependent perturbation,\\
mixed-exponent recovery, and the classical even-power background}
\author{C. Wayne Baker}
\date{Version 2.4 draft for author review --- November 10, 2026\\ Original manuscript: June 3, 2026}

\begin{document}
\maketitle

\begin{abstract}
Regular refinement families for spherical area exhibit the extrapolation ladder of orders
$2,4,6,8$. This manuscript asks where that ladder ends. Our main result is a cubic
obstruction: when mesh nodes are displaced by an amount proportional to the grid spacing, the
even-power error expansion acquires an odd $h^3$ term with coefficient $\varepsilon C_2'(0)$,
which the even-power operators cannot remove. They leave it with exact residual factors
$-1/6$, $-1/90$ and $-1/810$ at two, three and four scales, so the extrapolated order stalls
at three. For uniform square-cell grids split along a fixed diagonal and deformed by a fixed
smooth, boundary-vanishing field, we prove the underlying even-power expansion uniformly in
the displacement, give the cell coefficient in closed form, and certify by interval
arithmetic that $C_2'(0)=-6.80779\ldots\times10^{-4}\neq0$ for the test field, so that for this
family the third-order stall is a theorem. A deterministic numerical test on a chordally triangulated latitude--longitude cell
confirms the prediction: the fitted cubic coefficient agrees with $\varepsilon C_2'(0)$,
measured independently from fixed-displacement runs, to about three significant figures; the
residual ratios approach the predicted $1/15$ and $1/9$; the extrapolated orders stall at
three (three- and four-scale: $3.000$); and a mixed-exponent operator built for the exponent set $\{2,3\}$ restores fourth
order. We give exact weights and first surviving moments for the mixed sets $\{2,3\}$,
$\{2,3,4\}$ and $\{2,3,4,5\}$, and an elementary parity proof of the even-power expansion for
spherical caps bounded by regular geodesic polygons. The ladder itself, for flat-triangle
(chordal) triangulations of a fixed smooth parametrization, is classical: it was conjectured
by Georg (1991) and proved by Verlinden and Cools (1994) and by Lyness (1994). We verify that
the chordal spherical triangle, fan-refined polygon and uniform latitude--longitude cell of
this study satisfy that theorem, give an independent direct proof for the latitude--longitude
cell, and note that a 1994 experiment of Georg and Tausch coincides with our octant test. The
scale weights are the classical Richardson--Romberg weights. The new proofs and the
computer-assisted certificate were produced with AI assistance and await independent human
review. Nothing is claimed for independent node jitter, adaptive meshes or arbitrary
spherical meshes, and none of the classical ingredients is claimed as new.
\end{abstract}

%------------------------------------------------------------
\section{Introduction}
Many geometric area approximations admit structured error expansions in a refinement
parameter $h\asymp N^{-1}$. When the errors are in even powers, the same approximation
evaluated at successive scales can be combined to remove those terms, giving the sequence of
orders
\[
2\longrightarrow 4\longrightarrow 6\longrightarrow 8 .
\]
For flat-triangle triangulations of a fixed smooth surface patch this structure is a theorem
of the boundary-element literature (Section~\ref{sec:classical}). The question studied here is
what happens at the edge of that theorem. In practice a mesh is often not the image of a
uniform grid under one fixed map: nodes are smoothed, jittered or adapted by amounts tied to
the local spacing. The perturbed-grid experiment of the original June 2026 manuscript showed
the higher-order ladder deteriorating in exactly that situation while second-order
convergence survived.

Section~\ref{sec:perturb} gives a mechanism. A displacement $\delta=\varepsilon h$ turns the
coefficient $C_2(\delta)$ of the even-power expansion into a function of $h$, and its first
variation appears as an $h^3$ term. Because the even-power operators are built to annihilate
$h^2,h^4,h^6$, they multiply an $h^3$ term by fixed nonzero factors, which we compute
exactly. Theorem~\ref{thm:joint} proves the required expansion for smoothly deformed
square-cell grids, and Proposition~\ref{prop:certified} certifies that the cubic coefficient
is nonzero for the test field. A deterministic test, specified completely and included in the archive, confirms the
amplitude law, the residual factors and the stall at third order.
Section~\ref{sec:general} gives the operators that remove a prescribed mixed set of
exponents, and the same test confirms that the $\{2,3\}$ operator restores fourth order.

The remaining sections record the fixed-map families: a spherical cap bounded by a regular
geodesic polygon, for which we give a parity proof (Section~\ref{sec:cap}); the uniform
latitude--longitude cell, with an independent direct proof (Section~\ref{sec:latlon}); and the
chordal spherical triangle and fan-refined polygon, which illustrate the classical theorem
(Sections~\ref{sec:triangle} and~\ref{sec:polygon}). Earlier versions of this manuscript were
written without knowledge of the work of Georg~\cite{georg91}, Georg and
Tausch~\cite{georgtausch}, Verlinden and Cools~\cite{verlindencools} and
Lyness~\cite{lyness94}, and presented the triangle and polygon ladders as numerical evidence
for a conjectural law. The original tabulated experiments are retained with their provenance
limitations.

%------------------------------------------------------------
\section{Contributions and scope}\label{sec:scope}
\paragraph{Results of this manuscript.}
\begin{itemize}
\item \emph{A cubic obstruction} (Proposition~\ref{prop:cubic}): under a joint even-power
expansion in the mesh size and a displacement parameter, scale-dependent displacement
produces an $h^3$ term $\varepsilon C_2'(0)h^3$, and the two-, three- and four-scale
even-power operators leave residuals $-\tfrac16$, $-\tfrac1{90}$, $-\tfrac1{810}$ times that
term.
\item \emph{A displacement-uniform expansion} (Theorem~\ref{thm:joint}): for uniform
square-cell grids with a fixed diagonal, deformed by a fixed smooth boundary-vanishing field,
the chordal area has an even-power expansion with smooth coefficients and a remainder uniform
in the displacement. This supplies the hypothesis of Proposition~\ref{prop:cubic} for that
family (Corollary~\ref{cor:cubicfamily}).
\item \emph{The cell coefficient in closed form} (Lemma~\ref{lem:q2}), verified by exact
polynomial algebra.
\item \emph{A certified nonzero cubic coefficient} (Proposition~\ref{prop:certified}): a
computer-assisted interval enclosure of $C_2'(0)$ for the test field.
\item \emph{Numerical confirmation} (Section~\ref{sec:numtest}): a deterministic test of the
amplitude law, the residual ratios and the third-order stall, with the fixed-displacement
family as a control.
\item \emph{Mixed-exponent operators} (Section~\ref{sec:general}): exact weights, first
surviving moments and amplification factors for the sets $\{2,3\}$, $\{2,3,4\}$,
$\{2,3,4,5\}$, with numerical confirmation that the $\{2,3\}$ operator restores fourth order.
\item \emph{The spherical cap family} (Section~\ref{sec:cap}): exact spherical triangle areas
over a regular geodesic boundary polygon, which is not a flat-triangle rule, with an
elementary parity proof of its even-power expansion.
\item \emph{An independent direct proof} for uniform latitude--longitude chordal cells
(Proposition~\ref{prop:latlon-even}); the statement itself also follows from the classical
theorem.
\item Reproducible high-precision tables for the spherical test families, and the observation
that the octant case reproduces a published 1994 tableau (Remark~\ref{rem:gt}).
\end{itemize}

\paragraph{Classical background.}
Richardson--Romberg extrapolation~\cite{richardson,brs}; the Euler--Maclaurin expansion for
quadrature over triangles and its extrapolation~\cite{lyness78}; the flat-triangle
approximation of curved surfaces, its extrapolation and its even-power
expansion~\cite{georg91,georgtausch,verlindencools,lyness94}. By
Corollaries~\ref{cor:triangle}--\ref{cor:latlon} below, the even-power structure of the
chordal spherical triangle, the fan-refined spherical polygon and the uniform
latitude--longitude cell belongs to this classical material.

\paragraph{Not claimed.}
No new extrapolation method: the mixed-exponent operators are Richardson extrapolation for a
general exponent sequence~\cite{bulirschstoer,sidi}. No priority for extrapolating
flat-triangle surface approximations. No claim that the degradation of high-order behaviour
under perturbations proportional to the mesh size is new (see below). No even-power theorem
for independent node jitter, adaptive meshes or arbitrary spherical meshes.

\paragraph{Related work on perturbed meshes.}
That extrapolation needs an asymptotic expansion, and that such expansions are tied to
uniform or smoothly mapped meshes, is classical~\cite{marchukshaidurov,blumlinrannacher}.
Node perturbations proportional to the mesh size are a standard device for constructing
irregular test grids in the verification of finite-volume schemes, where grids that cannot be
smoothly mapped to a regular one are treated as a separate class~\cite{diskinthomas}; the
trapezoidal rule with nodes displaced by a fraction of the spacing has been analysed for
convergence~\cite{austintrefethen}. We did not find, in a limited search, the specific
statements proved here for chordal area: the displacement-uniform expansion, the closed-form
cell coefficient, and the exact cubic term with a certified nonzero coefficient. They are
offered as new to the author's knowledge only; given the joint expansion, the cubic term
itself is a one-line Taylor step.

%------------------------------------------------------------
\section{N-Series scale operators}\label{sec:operators}
Put $h=N^{-1}$, let $TA(h)=A(h/2)$, and set $A_N=A(h)$. If
$A(h)=A+c_2h^2+c_4h^4+c_6h^6+O(h^8)$, define
\begin{align}
\Ahat^{(2)}_N &= \frac{4A_{2N}-A_N}{3}, \label{eq:R2}\\
\Ahat^{(3)}_N &= \frac{A_N}{45}-\frac{4A_{2N}}{9}+\frac{64A_{4N}}{45}, \label{eq:R3}\\
\Ahat^{(4)}_N &= -\frac{A_N}{2835}+\frac{4A_{2N}}{135}-\frac{64A_{4N}}{135}+\frac{4096A_{8N}}{2835}. \label{eq:R4}
\end{align}
Each set of weights sums to one and eliminates $h^2$ and, respectively, the additional powers
$h^4$ and $h^6$. Hence, subject to a sufficiently long asymptotic expansion and nonzero next
coefficient, the errors are of orders 4, 6, and 8. These are Richardson--Romberg extrapolation
formulas in a geometric N-Series setting.

In terms of the shift $T$ (which acts on a pure power $h^p$ as multiplication by $2^{-p}$),
Eqs.~\eqref{eq:R2}--\eqref{eq:R4} are the nested Romberg stages
\[
\Ahat^{(r)}=\prod_{k=1}^{r-1}\frac{4^{k}T-I}{4^{k}-1}\,A,\qquad r=2,3,4 .
\]
For example, $(16T-I)(4T-I)/45=(64T^2-20T+I)/45$ reproduces Eq.~\eqref{eq:R3}. On an $h^p$
term, stage $k$ contributes the factor $(4^k2^{-p}-1)/(4^k-1)$, which vanishes exactly when
$p=2k$. This product form gives the residual moments used in Section~\ref{sec:perturb}.

%------------------------------------------------------------
\section{The regularity hypothesis}\label{sec:regularity}
For a fixed spherical domain $D$ and a specified geometrical area rule, call a refinement
family \emph{even-asymptotic through order $2r$} when, as $h\downarrow0$ along the admissible
refinement sequence,
\begin{equation}\label{eq:evenasym}
A(h)-A(D)=\sum_{k=1}^{r}c_{2k}h^{2k}+O(h^{2r+2}),
\end{equation}
with coefficients independent of the refinement level. Equation~\eqref{eq:evenasym}, rather
than informal smoothness alone, is the operative hypothesis of the $2,4,6,\dots$ ladder.

For the flat-triangle area rule, Theorem~\ref{thm:classical} below identifies a class for which
\eqref{eq:evenasym} holds to every order: a parametrization that is fixed (independent of
$h$), smooth and regular on a closed triangle, sampled on the uniform subdivision of that
triangle. Outside that class, and for other area rules, even-power structure must still be
proved for the particular geometry and rule or supported by convergence tests. The map should
preserve orientation and describe the same target domain at every scale.

%------------------------------------------------------------
\section{Background: the classical expansion for chordal triangulations}\label{sec:classical}
This section records known results and verifies their hypotheses for the constructions
used later. Nothing in it is claimed as new except the routine verifications.
Let $\sigma=\{(u,v):u\ge0,\ v\ge0,\ u+v\le1\}$. For $n\ge1$ the \emph{regular subdivision}
of $\sigma$ consists of the $n(n+1)/2$ triangles with vertices
$\bigl(\tfrac in,\tfrac jn\bigr),\bigl(\tfrac{i+1}n,\tfrac jn\bigr),\bigl(\tfrac in,\tfrac{j+1}n\bigr)$,
$0\le i+j\le n-1$, together with the $n(n-1)/2$ triangles with vertices
$\bigl(\tfrac{i+1}n,\tfrac jn\bigr),\bigl(\tfrac{i+1}n,\tfrac{j+1}n\bigr),\bigl(\tfrac in,\tfrac{j+1}n\bigr)$,
$0\le i+j\le n-2$: in all, $n^2$ congruent subtriangles. For a map $\phi:\sigma\to\R^3$ let
$S_n(\phi)$ denote the sum, over these $n^2$ subtriangles, of the Euclidean area of the flat
triangle whose vertices are the images under $\phi$ of the subtriangle's vertices.

\begin{theorem}[Verlinden--Cools~\cite{verlindencools}; Lyness~\cite{lyness94}]
\label{thm:classical}
Let $\phi:\sigma\to\R^3$ have continuous partial derivatives of all orders up to the boundary,
with $\phi_u\times\phi_v\neq0$ on $\sigma$. Then there are constants $c_1,c_2,\dots$ such that
\begin{equation}\label{eq:classical}
S_n(\phi)\ \sim\ \iint_\sigma\|\phi_u\times\phi_v\|\,du\,dv+\sum_{k\ge1}\frac{c_k}{n^{2k}}
\qquad(n\to\infty),
\end{equation}
the leading term being the area of the surface $\phi(\sigma)$.
\end{theorem}

This is the case $g\equiv1$ of Theorem~1 of~\cite{verlindencools}, whose operator weights each
flat-triangle area by a smooth function $g$ and which also yields the corresponding
expansions for Georg's modified trapezoidal and midpoint rules for $\int_S f\,dS$. It is
Conjecture~5.1 of~\cite{georg91}, restated as Conjecture~1 of~\cite{georgtausch}.
Lyness~\cite{lyness94} obtains the same $h^2$-expansion (his Theorem~5.6) with finite
smoothness and a correspondingly truncated expansion. The expansion is asymptotic; no
convergence of the infinite series is asserted. The statement requires $\phi$ to be the same
map at every $n$.

\begin{lemma}[Radial projection of a flat triangle]\label{lem:radial}
Let $A,B,C\in\R^3$ be linearly independent, put $P(u,v)=(1-u-v)A+uB+vC$ and
$\phi=P/\|P\|$. Then $P\neq0$ on $\sigma$, $\phi$ is real-analytic on a neighbourhood of
$\sigma$, and
\begin{equation}\label{eq:jac}
\phi_u\times\phi_v=\frac{\det[A,B,C]}{\|P\|^{4}}\,P,\qquad
\|\phi_u\times\phi_v\|=\frac{|\det[A,B,C]|}{\|P\|^{3}}\neq0 .
\end{equation}
\end{lemma}
\begin{proof}
$P(u,v)$ is a combination of $A,B,C$ with coefficients summing to one, so it cannot vanish
when they are independent; hence $\phi$ is analytic near $\sigma$. Write $a=B-A$, $b=C-A$ and
let $\Pi=I-\phi\phi^{\mathsf T}$. Then $\phi_u=\Pi a/\|P\|$ and $\phi_v=\Pi b/\|P\|$. For
vectors projected onto the plane orthogonal to the unit vector $\phi$,
$\Pi a\times\Pi b=\bigl(\phi\cdot(a\times b)\bigr)\phi$, and
$P\cdot(a\times b)=A\cdot\bigl((B-A)\times(C-A)\bigr)=\det[A,B,C]$ because the terms in $u$ and
$v$ vanish. Substituting gives \eqref{eq:jac}.
\end{proof}

When $A,B,C$ are unit vectors, the radial projection carries the edges of the flat triangle
to great-circle arcs, so $\phi(\sigma)$ is the spherical triangle $ABC$, and the map
$(i,j)\mapsto\phi(i/N,j/N)$ is exactly the normalized-barycentric node set used in the
experiments.

\begin{corollary}[Chordal spherical triangle]\label{cor:triangle}
Let $A,B,C$ be linearly independent unit vectors and let $A_N$ be the sum of the Euclidean
areas of the $N^2$ chordal triangles of the normalized-barycentric triangulation of the
spherical triangle $ABC$. Then, with $h=1/N$, the family is even-asymptotic to every order:
\[
A_N\ \sim\ \mathrm{Area}(ABC)+\sum_{k\ge1}c_{k}h^{2k}.
\]
\end{corollary}
\begin{proof}
Lemma~\ref{lem:radial} shows that $\phi=P/\|P\|$ satisfies the hypotheses of
Theorem~\ref{thm:classical}, and $A_N=S_N(\phi)$.
\end{proof}

\begin{corollary}[Fan-refined spherical polygon]\label{cor:fan}
Let $v_0,\dots,v_{m-1}$ be unit vectors such that each triple $(v_0,v_i,v_{i+1})$,
$1\le i\le m-2$, is linearly independent, and let $A_N$ be the sum over $i$ of the chordal
areas of Corollary~\ref{cor:triangle} for the triangles $v_0v_iv_{i+1}$. Then $A_N$ has an
asymptotic expansion in even powers of $h=1/N$ whose constant term is the sum of the areas of
the spherical triangles $v_0v_iv_{i+1}$. That constant equals the area of the polygonal domain
$D$ exactly when those triangles have disjoint interiors and union $D$.
\end{corollary}
\begin{proof}
A finite sum of expansions of the form \eqref{eq:classical}.
\end{proof}

\begin{corollary}[Uniform latitude--longitude cell]\label{cor:latlon}
Fix $-\pi/2<\phi_0<\phi_1<\pi/2$ and $\lambda_0<\lambda_1$. Subdivide the parameter rectangle
into a uniform $N\times N$ grid, split every cell along its $(p_{00},p_{11})$ diagonal, and let
$A_N$ be the total Euclidean area of the resulting chordal triangles. Then $A_N$ has an
asymptotic expansion in even powers of $h=1/N$ with constant term
$(\lambda_1-\lambda_0)(\sin\phi_1-\sin\phi_0)$.
\end{corollary}
\begin{proof}
Cut the rectangle along its own $(p_{00},p_{11})$ diagonal into two triangles. The uniform
grid with consistently oriented cell diagonals is the image of the regular subdivision of
$\sigma$ under the affine parametrization of each of the two triangles: the first family of
subtriangles maps to the cell triangles $(p_{00},p_{10},p_{11})$ and the second to the cell
triangles $(p_{00},p_{11},p_{01})$ for one large triangle, and conversely for the other.
Composing with the spherical coordinate map, which is analytic with nonvanishing normal
$\cos\phi\neq0$ on the closed band, gives two maps satisfying Theorem~\ref{thm:classical}.
\end{proof}

\begin{remark}[The 1994 experiment of Georg and Tausch]\label{rem:gt}
Section~3 of~\cite{georgtausch} integrates $f=x^2+y^2+z^2$ over the positive octant of the
unit sphere, parametrized by radial projection of the triangle with vertices
$ae_1,be_2,ce_3$, on meshes $h_i=2^{-i}$, and forms a Romberg tableau. Since $f=1$ at every
mapped vertex, their modified trapezoidal rule with $a=b=c=1$ is the chordal area $A_N$ of the
octant, which is the $\beta=90^\circ$ case of Section~\ref{sec:triangle}.
Table~\ref{tab:gt} confirms that the two computations agree to the digits printed in 1994.
Their tableau also shows the $h^2,h^4,h^6$ columns. The octant experiment of this manuscript
is therefore a higher-precision repetition of a published one.
\end{remark}

\begin{table}[ht]
\centering
\begin{tabular}{rll}
\toprule
$N$ & $|A_N-\pi/2|$, this archive & Georg--Tausch (1994), trapezoidal, $a=b=c=1$\\
\midrule
4 & $7.75\times10^{-2}$ & $7.7\times10^{-2}$\\
8 & $2.03\times10^{-2}$ & $2.0\times10^{-2}$\\
16 & $5.14\times10^{-3}$ & $5.1\times10^{-3}$\\
32 & $1.29\times10^{-3}$ & $1.3\times10^{-3}$\\
\bottomrule
\end{tabular}
\caption{Raw chordal-area error for the octant. Left: recomputed on November~10, 2026 from the
archived routine at 50 digits. Right: first column of the tableau printed
in~\cite{georgtausch}.}\label{tab:gt}
\end{table}

%------------------------------------------------------------
\section{Main result: scale-dependent perturbations and the cubic obstruction}\label{sec:perturb}
In the original perturbation experiment, boundary nodes of a latitude--longitude cell
(Section~\ref{sec:latlon}) remained
fixed, while interior nodes were displaced by a deterministic smooth field. Perturbation
amplitude was specified as a fraction of local grid spacing; three seeds were used at each
level. Because the amplitude is a fraction of the local spacing, the parametrization changes
with $N$ and Theorem~\ref{thm:classical} does not apply.

\begin{table}[ht]
\centering
\begin{tabular}{lrrrr}
\toprule
Relative perturbation & Raw & Two-scale & Three-scale & Four-scale\\
\midrule
0.00 & 1.99996836 & 3.99999637 & 5.99994223 & 7.99911861\\
0.02 & 1.99997529 & 4.02004815 & 4.53899093 & 2.61899119\\
0.05 & 2.00001207 & 3.44420303 & 5.68235637 & 2.77050496\\
0.10 & 2.00014360 & 3.89815337 & 8.94059560 & 3.59334301\\
0.20 & 2.00066838 & 3.90830154 & 6.15480240 & 4.36890149\\
\bottomrule
\end{tabular}
\caption{Original perturbed-grid measurements (Table~6 of the June 2026 manuscript):
aggregate final-level mean observed orders. These are diagnostics, not verified asymptotic
exponents.}\label{tab:perturb}
\end{table}

An apparent order such as $8.94059560$ is not evidence of asymptotic convergence above the
theoretical order: accidental coefficient cancellation, small denominators in measured error
ratios, sign changes, or numerical floors can distort finite-level slopes. The original data
establish only that the second-order raw behavior is comparatively robust and that
higher-order extrapolation is not uniform in this test window.

\paragraph{Qualitative reading of Table~\ref{tab:perturb}.}
At small perturbation the two-scale order stays near 4 (0.02: 4.02) while the three- and
four-scale orders fall well below their nominal 6 and 8, with four-scale values between about
2.6 and 4.4. This is qualitatively what a residual $h^3$ term would produce: the higher-order
operators remove the $h^4,h^6$ terms that otherwise mask it, and the residual factors of
Proposition~\ref{prop:cubic} below (e.g.\ $-1/6$ versus $-1/810$) are applied to a coefficient
that is linear in the amplitude. This is compatibility, not confirmation. The non-monotone
entries, the erratic three-scale values, and the demonstrable pre-asymptotic variation in the
unperturbed four-scale baseline (cf.\ Table~\ref{tab:cap}, $30^\circ$) mean that the table
cannot by itself separate a cubic residual from finite-window or round-off effects.

\subsection{Conditional appearance of a cubic term}
The key distinction is between a fixed nonuniform parametrization and a parametrization whose
displacement amplitude changes with $h$. Suppose $A(h,\delta)$ is an approximation defined
using a perturbation parameter $\delta$, and the following smooth, uniform in $\delta$
even-power expansion holds near $(h,\delta)=(0,0)$:
\begin{equation}\label{eq:joint}
A(h,\delta)-A=C_2(\delta)h^2+C_4(\delta)h^4+O(h^6),
\end{equation}
where $C_2$ is three times continuously differentiable (or has a locally Lipschitz second
derivative) and $C_4$ is at least once continuously differentiable near zero.

\begin{remark}[Status of the hypothesis]\label{rem:supplies}
Let the displacement be $\delta V$ for a fixed smooth field $V$ on the parameter domain that
vanishes on its boundary. For each fixed small $\delta$, the displaced grid is the image of
the uniform grid under the fixed diffeomorphism $\psi_\delta=\mathrm{id}+\delta V$, so
Theorem~\ref{thm:classical} already gives an even-power expansion with the same constant term
$A$. Theorem~\ref{thm:joint} below proves more: for uniform square-cell grids with a fixed
diagonal the coefficients are $C^\infty$ in $\delta$ and the remainder is uniform, so
\eqref{eq:joint} holds in full for that family. The local cubic terms that must cancel
in~\cite{georgtausch} cancel only by summation over a uniform subdivision of a fixed map,
which is consistent with a cubic term reappearing when the map itself varies with $h$.
\end{remark}

\begin{proposition}[Scale-dependent perturbations can create odd powers]\label{prop:cubic}
Assume \eqref{eq:joint} with a remainder uniform for $|\delta|\le\delta_0$ and set
$\delta=\varepsilon h$, with fixed $\varepsilon$. Then
\begin{equation}\label{eq:cubic}
A(h,\varepsilon h)-A=C_2(0)h^2+\varepsilon C_2'(0)h^3
+\Big[\tfrac12\varepsilon^2C_2''(0)+C_4(0)\Big]h^4+O(h^5).
\end{equation}
If $\varepsilon C_2'(0)\neq0$, then the two-, three- and four-scale even-power operators
\eqref{eq:R2}--\eqref{eq:R4} each leave a cubic residual
\[
-\tfrac16\,\varepsilon C_2'(0)\,h^3,\qquad
-\tfrac1{90}\,\varepsilon C_2'(0)\,h^3,\qquad
-\tfrac1{810}\,\varepsilon C_2'(0)\,h^3,
\]
respectively, so the extrapolated error is $O(h^3)$ and is not improved beyond third order by
adding scales.
\end{proposition}

\begin{proof}
Taylor-expand $C_2(\varepsilon h)$ and $C_4(\varepsilon h)$ in \eqref{eq:joint}; the $h^3$
coefficient is $\varepsilon C_2'(0)$. For a scale-doubling weight set $\{w_j\}$ the surviving
coefficient is $\varepsilon C_2'(0)\sum_jw_j2^{-3j}$. Direct summation gives
\begin{align*}
\text{(2)}&:\ -\tfrac13+\tfrac43\cdot\tfrac18=-\tfrac16,\\
\text{(3)}&:\ \tfrac1{45}-\tfrac49\cdot\tfrac18+\tfrac{64}{45}\cdot\tfrac1{64}=-\tfrac1{90},\\
\text{(4)}&:\ -\tfrac1{2835}+\tfrac4{135}\cdot\tfrac18-\tfrac{64}{135}\cdot\tfrac1{64}
+\tfrac{4096}{2835}\cdot\tfrac1{512}=-\tfrac1{810}.
\end{align*}
Equivalently, by the product form in Section~\ref{sec:operators} the stage factors at $p=3$
are $-\tfrac16$, $\tfrac1{15}$, $\tfrac19$, whose cumulative products are
$-\tfrac16,-\tfrac1{90},-\tfrac1{810}$.
\end{proof}

\begin{remark}
The notation $A(h,\varepsilon h)$ alone does not establish the existence of a cubic error
term. It might vanish if $C_2'(0)=0$, if symmetry forces cancellations, or if a different
mesh/area discretization fails the joint expansion assumption. Moreover, a fixed physical
displacement $\delta$ and an $h$-scaled displacement define different refinement families; only
comparisons made with identical geometric rules and numerical precision isolate this effect.
\end{remark}

\subsection{The joint expansion as a theorem}\label{sec:jointthm}
\begin{theorem}[Displacement-uniform even-power expansion]\label{thm:joint}
Let $\Omega=[0,1]^2$, let $X$ be a $C^\infty$ map from a neighbourhood of $\Omega$ into $\R^3$
with $X_s\times X_t\neq0$ on $\Omega$, and let $V$ be a $C^\infty$ planar vector field on a
neighbourhood of $\Omega$ that vanishes on $\partial\Omega$. Let $\delta_0>0$ be small enough
that $\delta_0\sup_\Omega\|DV\|<\tfrac12$, and put $\psi_\delta=\mathrm{id}+\delta V$ and
$F_\delta=X\circ\psi_\delta$ for $|\delta|\le\delta_0$. For $N\ge1$ and $h=1/N$, sample
$F_\delta$ on the uniform $(N+1)\times(N+1)$ parameter grid, split every cell along its
$(p_{00},p_{11})$ diagonal, and let $A(h,\delta)$ be the sum of the Euclidean areas of the
resulting chordal triangles. Let $A=\iint_\Omega\|X_s\times X_t\|\,ds\,dt$. Then for every
$r\ge1$ there are $C^\infty$ functions $C_2,\dots,C_{2r}$ on $[-\delta_0,\delta_0]$ and a
constant $M_r$ such that, for all sufficiently large $N$,
\begin{equation}\label{eq:jointthm}
\sup_{|\delta|\le\delta_0}\Bigl|A(h,\delta)-A-\sum_{k=1}^{r}C_{2k}(\delta)h^{2k}\Bigr|
\le M_rh^{2r+2}.
\end{equation}
Moreover, with $q_0,q_2$ the cell coefficients of Lemma~\ref{lem:q2} for $F=F_\delta$,
\begin{equation}\label{eq:C2formula}
C_2(\delta)=\iint_\Omega q_2(x,\delta)\,dx-\frac1{24}\iint_\Omega\Delta_xq_0(x,\delta)\,dx .
\end{equation}
\end{theorem}
\begin{proof}
\emph{Step 1: a uniform family of regular maps.} For $|\delta|\le\delta_0$ the estimate
$\|\psi_\delta(x)-\psi_\delta(y)\|\ge(1-|\delta|\sup\|DV\|)\|x-y\|$ shows that $\psi_\delta$
is injective, and $\det D\psi_\delta>0$. Since $V=0$ on $\partial\Omega$, $\psi_\delta$ fixes
$\partial\Omega$ pointwise, so by injectivity no interior point is mapped to $\partial\Omega$.
For an interior point $x_0$ the path $\delta'\mapsto\psi_{\delta'}(x_0)$, $\delta'$ between $0$
and $\delta$, starts at $x_0$ and never meets $\partial\Omega$; hence $\psi_\delta$ maps the
interior of $\Omega$ into itself. The image of the interior is open, because $\psi_\delta$ is a local diffeomorphism. It is
also relatively closed in the interior: if $y_n=\psi_\delta(x_n)$ with $x_n$ interior and
$y_n\to y$ for an interior point $y$, then by compactness of $\Omega$ a subsequence of $(x_n)$
converges to some $x\in\Omega$, and $\psi_\delta(x)=y$ by continuity; $x$ cannot lie on
$\partial\Omega$, since then $\psi_\delta(x)=x\in\partial\Omega$ whereas $y$ is interior, so $y$
is the image of an interior point. The interior of $\Omega$ is connected, so a nonempty subset
that is both open and relatively closed is the whole interior. Thus
$\psi_\delta$ is a bijection of $\Omega$, and the change of variables formula gives
$\iint_\Omega\|(F_\delta)_s\times(F_\delta)_t\|=\iint_\Omega\|X_s\times X_t\|\circ\psi_\delta\,
\det D\psi_\delta=A$ for every $\delta$. The family $F_\delta$ is jointly $C^\infty$ in
$(x,\delta)$ on a neighbourhood of $\Omega\times[-\delta_0,\delta_0]$, and
$\|(F_\delta)_s\times(F_\delta)_t\|$ is bounded below by a positive constant there.

\emph{Step 2: the paired cell is even in $h$.} For a cell centre $x$ put
$p_{ab}(x,h)=F_\delta\bigl(x+((2a-1)h/2,(2b-1)h/2)\bigr)$, $a,b\in\{0,1\}$, and
\[
G(x,h,\delta)=\tfrac12\|(p_{10}-p_{00})\times(p_{11}-p_{00})\|
+\tfrac12\|(p_{11}-p_{00})\times(p_{01}-p_{00})\| .
\]
Each difference of corners is $h$ times a jointly smooth vector function, so each cross
product is $h^2w_i(x,h,\delta)$ with $w_i$ jointly smooth and
$w_i(x,0,\delta)=(F_\delta)_s\times(F_\delta)_t\neq0$. By compactness there is $h_0>0$ such
that $w_i\neq0$ for $|h|\le h_0$, so $g=G/h^2=\tfrac12(\|w_1\|+\|w_2\|)$ is jointly
$C^\infty$ on $\Omega\times[-h_0,h_0]\times[-\delta_0,\delta_0]$, the smooth extensions of
$X$ and $V$ being used where a shifted corner lies outside $\Omega$. Replacing $h$ by $-h$
interchanges $p_{00}\leftrightarrow p_{11}$ and $p_{10}\leftrightarrow p_{01}$, which
exchanges the two unoriented triangles; hence $g(x,-h,\delta)=g(x,h,\delta)$. Taylor's
theorem in $h$ therefore gives
\[
G(x,h,\delta)=h^2\Bigl(q_0(x,\delta)+q_2(x,\delta)h^2+\dots+q_{2r}(x,\delta)h^{2r}\Bigr)
+O(h^{2r+4}),
\]
with jointly $C^\infty$ coefficients, $q_0=\|(F_\delta)_s\times(F_\delta)_t\|$, and a remainder
uniform in $(x,\delta)$.

\emph{Step 3: midpoint summation.} $A(h,\delta)$ is the sum of $G$ over the $N^2$ cell centres
$x_{ij}=((i+\tfrac12)h,(j+\tfrac12)h)$, so the remainders of Step~2 contribute $O(h^{2r+2})$
uniformly. For a smooth function $q$ on $\Omega$, applying the one-dimensional midpoint
Euler--Maclaurin formula in each variable (see~\cite{davisrabinowitz}, or Lemma~1
of~\cite{verlindencools} with offset $\tfrac12$) gives
\[
h^2\sum_{i,j=0}^{N-1}q(x_{ij})=\iint_\Omega q\,dx-\frac{h^2}{24}\iint_\Omega\Delta q\,dx
+\sum_{\ell=2}^{r}h^{2\ell}B_{2\ell}[q]+O(h^{2r+2}),
\]
where each $B_{2\ell}[q]$ is the integral over $\Omega$ of a fixed linear differential
expression in $q$, and the remainder is bounded by a constant times a finite number of
sup-norms of derivatives of $q$. Applied to $q=q_{2k}(\cdot,\delta)$, the coefficients are
$C^\infty$ in $\delta$ (differentiation under the integral) and the remainders are uniform in
$\delta$ (compactness). Collecting powers of $h$ gives \eqref{eq:jointthm}, the constant term
$\iint q_0=A$ by Step~1, and \eqref{eq:C2formula}.
\end{proof}

\begin{lemma}[Cell coefficient]\label{lem:q2}
Let $F$ be a $C^\infty$ map into $\R^3$ near $x$ with $N=F_s\times F_t\neq0$ at $x$, and let
$G(x,h)$ be the paired-cell chordal area of Step~2 above. Then
$G=h^2\bigl(q_0+q_2h^2+O(h^4)\bigr)$ with
\begin{equation}\label{eq:q2}
q_0=\|N\|,\qquad
q_2=\frac{R\cdot N}{48\,\|N\|}+\frac{(F_{st}\cdot N)^2\,\|F_s+F_t\|^2}{8\,\|N\|^3},
\end{equation}
where $R=(F_s-F_t)\times(\partial_s+\partial_t)^3F+(\partial_s-\partial_t)^3F\times(F_s+F_t)$
and all derivatives are evaluated at $x$.
\end{lemma}
\begin{proof}
Write $D_1=\partial_s+\partial_t$, $D_2=\partial_s-\partial_t$, and put
$c=p_{11}-p_{00}$, $e=p_{10}-p_{01}$, $m=\tfrac12(p_{10}+p_{01})-\tfrac12(p_{00}+p_{11})$.
Then $p_{10}-p_{00}=m+\tfrac12e+\tfrac12c$ and $p_{01}-p_{00}=m-\tfrac12e+\tfrac12c$, so the
two cross products are $(\tfrac12e+m)\times c$ and $(\tfrac12e-m)\times c$, and
\[
G=\tfrac12\|n+\mu\|+\tfrac12\|n-\mu\|,\qquad n=\tfrac12\,e\times c,\quad\mu=m\times c .
\]
Taylor expansion along the two diagonals gives $c=hD_1F+\tfrac{h^3}{24}D_1^3F+O(h^5)$,
$e=hD_2F+\tfrac{h^3}{24}D_2^3F+O(h^5)$ and
$m=\tfrac{h^2}{8}(D_2^2-D_1^2)F+O(h^4)=-\tfrac{h^2}{2}F_{st}+O(h^4)$. Since
$D_2F\times D_1F=2N$,
\[
n=h^2N+\tfrac{h^4}{48}R+O(h^6),\qquad\mu=-\tfrac{h^3}{2}\,F_{st}\times(F_s+F_t)+O(h^5).
\]
For vectors with $\|\mu\|<\|n\|$,
$\tfrac12(\|n+\mu\|+\|n-\mu\|)=\|n\|+\bigl(\|\mu\|^2-(\mu\cdot\hat n)^2\bigr)/(2\|n\|)
+O(\|\mu\|^4/\|n\|^3)$ with $\hat n=n/\|n\|$. The first term is
$h^2\|N\|+\tfrac{h^4}{48}R\cdot\hat N+O(h^6)$, $\hat N=N/\|N\|$. For the second, decompose
$F_{st}=u_T+\gamma\hat N$ with $u_T$ in the span of $F_s,F_t$ and $\gamma=F_{st}\cdot\hat N$.
Then $u_T\times(F_s+F_t)$ is parallel to $\hat N$, while $\hat N\times(F_s+F_t)$ is
orthogonal to $\hat N$ with norm $\|F_s+F_t\|$, so the part of $\mu$ orthogonal to $\hat N$
has squared norm $\tfrac{h^6}{4}\gamma^2\|F_s+F_t\|^2+O(h^8)$. Dividing by $2\|n\|=2h^2\|N\|+O(h^4)$
gives the second term of \eqref{eq:q2}.
\end{proof}
{\sloppy The identity \eqref{eq:q2} was also verified as an exact polynomial identity, with the
Taylor coefficients of $F$ through fourth order treated as 42 independent symbols
(\texttt{scripts/q2\_lemma\_symbolic\_audit.py}): the two leading normals equal $N$, the
$h^3$ terms of the two triangles cancel, the $h^4$ coefficient equals $q_2$, and the
fourth-order coefficients do not enter.\par}

\begin{corollary}[Cubic obstruction for smoothly deformed square-cell grids]
\label{cor:cubicfamily}
In the setting of Theorem~\ref{thm:joint}, fix $\varepsilon\in\R$ and take $\delta=\varepsilon h$.
Then \eqref{eq:cubic} holds, and the two-, three- and four-scale even-power operators leave the
cubic residuals $-\tfrac16$, $-\tfrac1{90}$, $-\tfrac1{810}$ times $\varepsilon C_2'(0)h^3$.
If $\varepsilon C_2'(0)\neq0$, each extrapolated error has exact asymptotic order three; if
$C_2'(0)=0$, no stall at third order follows.
\end{corollary}
\begin{proof}
Theorem~\ref{thm:joint} with $r=2$ gives \eqref{eq:joint} with $C^\infty$ coefficients and a
remainder uniform for $|\delta|\le\delta_0$, and $|\varepsilon h|\le\delta_0$ for large $N$.
Apply Proposition~\ref{prop:cubic}.
\end{proof}

\begin{proposition}[Certified cubic coefficient for the test field; computer-assisted]
\label{prop:certified}
Let $X$ be the unit-sphere latitude--longitude map of the cell
$30^\circ\le\phi\le60^\circ$, $0^\circ\le\lambda\le30^\circ$ in normalized coordinates
$(s,t)\in[0,1]^2$, and let
$V(s,t)=\sin(\pi s)\sin(\pi t)\bigl(\cos(2t+\tfrac3{10}),\sin(3s+\tfrac7{10})\bigr)$. Then
\begin{align*}
C_2(0)&\in-0.01189423662103184634603051\pm8.5\times10^{-27},\\
C_2'(0)&\in-6.80779112122198\times10^{-4}\pm2.0\times10^{-19}.
\end{align*}
In particular $C_2'(0)<0$, and by Corollary~\ref{cor:cubicfamily} every $\varepsilon\neq0$
gives a family whose even-power extrapolations have exact order three.
\end{proposition}
\begin{proof}
The four partial derivatives of $V$ are bounded on $\Omega$ by $\pi$, $\sqrt{\pi^2+4}$,
$\sqrt{\pi^2+9}$ and $\pi$, so the operator norm of $DV$ is at most its Frobenius norm
$\sqrt{4\pi^2+13}<7.25$. Hence $\delta_0=0.068$ is admissible in Theorem~\ref{thm:joint}, since
$0.068\times7.25=0.493<\tfrac12$. This covers every run of Section~\ref{sec:numtest}: the fixed
displacements satisfy $|\delta|\le0.04$ and the scale-dependent ones
$|\delta|=|\varepsilon|/N\le0.2/8=0.025$. Differentiate \eqref{eq:C2formula} under the
integral at $\delta=0$, using $F_\delta=X+\delta(X_sV_1+X_tV_2)+O(\delta^2)$ in \eqref{eq:q2}
and $\partial_\delta q_0|_{\delta=0}=\operatorname{div}(\|X_s\times X_t\|\,V)$. The resulting
integrand is an explicit combination of trigonometric functions whose only denominators are
powers of $\|X_s\times X_t\|=(\pi/6)^2\cos\phi$, with $\cos\phi\ge\tfrac12$ on the cell; it is
analytic on a complex neighbourhood of $\Omega$. The symbolic differentiation was carried out
in SymPy and the double integral enclosed by nested certified integration in Arb ball
arithmetic~\cite{johansson} (\texttt{scripts/c2prime\_enclosure.py}, 128-bit working
precision).
\end{proof}

\begin{remark}[What the certificate rests on]
\sloppy Proposition~\ref{prop:certified} is a computer-assisted statement: it relies on
Theorem~\ref{thm:joint}, Lemma~\ref{lem:q2}, and the correctness of SymPy's differentiation
and of Arb's enclosures. Three non-rigorous cross-checks agree with it: tensor
Gauss--Legendre quadrature of the same integrand in binary64 arithmetic reproduces both
midpoints to about $10^{-17}$ (\texttt{scripts/c2prime\_crosscheck\_numpy.py}); formula
\eqref{eq:C2formula}, evaluated from single cells without any grid sum, reproduces the
grid-sum values of $C_2(\delta)$ at $\delta=0,\pm0.02,\pm0.04$ to about $10^{-13}$
(\texttt{scripts/check\_joint\_expansion.py}); and direct finite differences of grid sums in
$\delta$ converge to the same value. The certificate concerns this field on this cell only.
\end{remark}

\subsection{Falsifiable cubic-coefficient test}
A sharper test of Proposition~\ref{prop:cubic} is the amplitude law, not merely a measured
convergence order. Choose a single smooth, boundary-vanishing perturbation vector field $V$
and a fixed chordal triangulation convention. Here $h=N^{-1}$ in normalized cell coordinates
(equivalently, $h=H/N$ with $H=1$); $\delta=\varepsilon/N$ is the dimensionless displacement
amplitude multiplying $V$. At signed amplitude $\varepsilon$, displace interior parameter-grid
nodes by $\delta V$; retain the same physical domain and check triangle orientation. With
$e(h,\varepsilon)=A(h,\varepsilon h)-A$, fit, over several successively refined grids and at
sufficiently high precision,
\begin{equation}\label{eq:fit}
e(h,\varepsilon)=b_2(\varepsilon)h^2+b_3(\varepsilon)h^3+b_4(\varepsilon)h^4+\cdots .
\end{equation}
Under the joint expansion \eqref{eq:joint}, the precise small-amplitude prediction is
$b_3(\varepsilon)=\varepsilon C_2'(0)$, including $b_3(0)=0$; a nonzero slope requires
$C_2'(0)\neq0$. Test symmetric amplitudes $0,\pm\varepsilon_1,\pm\varepsilon_2$ with at least
two distinct nonzero magnitudes. Check that
$[b_3(\varepsilon)-b_3(-\varepsilon)]/(2\varepsilon)$ stabilizes to one value, while the
even-in-$\varepsilon$ part of $b_3$ tends to zero; compare the inferred slope across
refinement windows and increased precision. By Remark~\ref{rem:supplies}, the fixed-$\delta$
runs of such a test should show clean even-power behaviour; that is a prediction of the
classical theorem and a useful control.

By Proposition~\ref{prop:cubic}, the predicted cubic coefficients of the extrapolated errors
are
\[
-\tfrac16\,\varepsilon C_2'(0)\ \ (\text{two-scale}),\qquad
-\tfrac1{90}\,\varepsilon C_2'(0)\ \ (\text{three-scale}),\qquad
-\tfrac1{810}\,\varepsilon C_2'(0)\ \ (\text{four-scale}).
\]
These coefficients are predictions of the proposition, not fitted measurements. A numerical test must
distinguish failure of the joint expansion, vanishing derivatives, and pre-asymptotic effects;
observed convergence orders alone cannot prove the hypothesis.

\subsection{Numerical confirmation on a specified deterministic test}\label{sec:numtest}
The test of the preceding subsection was carried out on November~10, 2026 with two separately
written archived scripts: \texttt{scripts/cubic\_amplitude\_test\_mpfr.py}, in 160-bit MPFR
arithmetic (about 48 significant digits, identical on every platform), and
\texttt{scripts/cubic\_amplitude\_test.py}, in the platform's extended precision. The cell is $30^\circ\le\phi\le60^\circ$,
$0^\circ\le\lambda\le30^\circ$, with a uniform $N\times N$ parameter grid in normalized
coordinates $(s,t)\in[0,1]^2$ and the $(p_{00},p_{11})$ diagonal. Nodes are displaced to
$(s,t)+\delta V(s,t)$ with the smooth, boundary-vanishing field
\[
V(s,t)=\sin(\pi s)\sin(\pi t)\,\bigl(\cos(2t+0.3),\ \sin(3s+0.7)\bigr).
\]
There are no random seeds. All numbers below are from the MPFR run. The extended-precision
run reproduces every one of them to the digits shown, except the three- and four-scale orders
of the fixed-displacement control, where extended precision reaches its rounding floor.
Boundary nodes are fixed by construction. Orientation was checked for each of the
$14{,}679{,}680$ triangles in the 67 runs: the signed area of the displaced triangle in the
parameter plane stayed positive, with minimum $0.878$ times its undisplaced value, and every
chordal normal pointed outward. No triangle failed either check.

\paragraph{Control: fixed displacement.}
For $\delta\in\{0,\pm0.02,\pm0.04\}$ held fixed as $N$ grows, the observed orders at
$N=32\to64$ are $2.0000$, between $4.0000$ and $4.0044$, between $5.997$ and $6.000$, and
between $7.998$ and $8.003$ for the raw, two-, three- and four-scale values. This is the even-power behaviour that
Theorems~\ref{thm:classical} and~\ref{thm:joint} predict for a fixed map. Extrapolating
$N^2(A_N-A)$ gives $C_2(0)=-0.011894236621$, in agreement with
Proposition~\ref{prop:certified} to all twelve digits. A two-step central difference of the
grid-sum values at $\delta=\pm0.02,\pm0.04$ gives $C_2'(0)\approx-6.806\times10^{-4}$; the
certified value is
\[
C_2'(0)=-6.80779\ldots\times10^{-4},
\]
the discrepancy in the fourth digit being the truncation error of that coarse difference.

\paragraph{Scale-dependent displacement.}
With $\delta=\varepsilon/N$, the cubic coefficient $b_3$ of \eqref{eq:fit} was estimated from
the two-scale residual as $-6N^3\bigl(\Ahat^{(2)}_N-A\bigr)$, with one extrapolation step in
$1/N$ between $N=64$ and $N=128$. Table~\ref{tab:b3} compares it with the prediction
$\varepsilon C_2'(0)$ from the certified coefficient. The two agree to about three
significant figures, the coefficient is odd in $\varepsilon$, and it doubles when
$\varepsilon$ doubles. At $N=64$, $\varepsilon=0.1$, the part of the two-scale residual that
is odd in $\varepsilon$ is $4.3269\times10^{-11}$, against the predicted
$-\tfrac16\varepsilon C_2'(0)h^3=4.328\times10^{-11}$.

\begin{table}[ht]
\centering
\begin{tabular}{rrr}
\toprule
$\varepsilon$ & fitted $b_3$ & predicted $\varepsilon C_2'(0)$\\
\midrule
$0.1$ & $-6.809\times10^{-5}$ & $-6.808\times10^{-5}$\\
$-0.1$ & $6.809\times10^{-5}$ & $6.808\times10^{-5}$\\
$0.2$ & $-1.3611\times10^{-4}$ & $-1.3616\times10^{-4}$\\
$-0.2$ & $1.3613\times10^{-4}$ & $1.3616\times10^{-4}$\\
\bottomrule
\end{tabular}
\caption{Amplitude law for the cubic coefficient (new in v2.4). Fitted values: 160-bit MPFR
grid sums. Predicted values: certified coefficient of Proposition~\ref{prop:certified}.}
\label{tab:b3}
\end{table}

\paragraph{Residual factors and the stall at third order.}
Proposition~\ref{prop:cubic} predicts that the three- and two-scale residuals are in the
ratio $(-1/90)/(-1/6)=1/15$ and the four- and three-scale residuals in the ratio
$(-1/810)/(-1/90)=1/9$. Table~\ref{tab:ratios} shows both ratios approaching those values,
and the extrapolated orders at or approaching three: adding scales no longer raises the
order, as the proposition states.

\begin{table}[ht]
\centering
\begin{tabular}{rrrrrrrr}
\toprule
& & \multicolumn{2}{c}{Residual ratios} & \multicolumn{4}{c}{Observed orders, $N\to2N$}\\
$\varepsilon$ & $N$ & three/two & four/three & Raw & Two & Three & Four\\
\midrule
$0.1$ & 64 & 0.069889 & 0.111099 & 2.0001 & 2.9656 & 3.0001 & 3.0000\\
$0.1$ & 128 & 0.068233 & 0.111108 & & & & \\
$0.2$ & 64 & 0.066837 & 0.111141 & 2.0001 & 2.9982 & 2.9997 & 3.0000\\
$0.2$ & 128 & 0.066767 & 0.111118 & & & & \\
\midrule
\multicolumn{2}{l}{Predicted} & 0.066667 & 0.111111 & 2 & 3 & 3 & 3\\
\bottomrule
\end{tabular}
\caption{Scale-dependent family $\delta=\varepsilon/N$: ratios of the extrapolated errors
and observed orders (new in v2.4; 160-bit MPFR).}\label{tab:ratios}
\end{table}

\paragraph{What this does and does not establish.}
By Corollary~\ref{cor:cubicfamily} and Proposition~\ref{prop:certified}, the third-order stall
for this cell and field is a theorem, subject to the review noted in the limitations; the
test above is its numerical illustration and an end-to-end check of the constants. The
fitted and certified cubic coefficients agree to three or four significant figures, the
remaining difference being the finite-$N$ error of the fit. The result concerns one field on
one cell. It does not reconstruct the original experiment of Table~\ref{tab:perturb}, whose
displacement field and seeds are not available, so the anomalies of that table are explained
in kind but not reproduced.

\subsection{Status of the original numerical comparison}
The previously reported numerical comparison of fixed-map and scale-dependent
perturbations has not been independently reproduced from the original
implementation and is omitted from this release; Section~\ref{sec:numtest} reports a new,
fully specified test instead. Proposition~\ref{prop:cubic} is an algebraic implication of
its stated joint expansion hypothesis; Theorem~\ref{thm:joint} verifies that hypothesis for
smooth fixed-field deformations of square-cell grids, not for the original perturbation code.

%------------------------------------------------------------
\section{Mixed-exponent operators: restoring the ladder}\label{sec:general}
Suppose the approximation has an expansion
\[
A(h)=A+\sum_{p\in P}c_ph^p+\text{higher-order terms},
\]
where $P=\{p_1,\dots,p_m\}$ is a prescribed set of distinct positive exponents. Using $m+1$
refinement levels $h,h/2,\dots,h/2^m$, seek
\begin{equation}\label{eq:gen}
\Ahat(h)=\sum_{j=0}^{m}w_jA(h/2^j),\qquad
\sum_{j=0}^{m}w_j=1,\qquad
\sum_{j=0}^{m}w_j2^{-jp_k}=0\quad(k=1,\dots,m).
\end{equation}
The associated Vandermonde system has a unique solution for distinct exponents; this is
Richardson extrapolation for a general exponent sequence~\cite{bulirschstoer,sidi}, and the
formulas below are worked instances of that classical process on a doubling ladder. If $P=\{2,3\}$, the three-scale
weights are
\begin{equation}\label{eq:P23}
\Ahat^{\{2,3\}}(h)=\frac{A(h)}{21}-\frac{4A(h/2)}{7}+\frac{32A(h/4)}{21}.
\end{equation}
For $P=\{2,3,4\}$ the four-scale formula is
\begin{equation}\label{eq:P234}
\Ahat^{\{2,3,4\}}(h)=-\frac{A(h)}{315}+\frac{4A(h/2)}{45}-\frac{32A(h/4)}{45}+\frac{512A(h/8)}{315}.
\end{equation}
For $P=\{2,3,4,5\}$ the five-scale formula is
\begin{equation}\label{eq:P2345}
\Ahat^{\{2,3,4,5\}}(h)=\frac{A(h)}{9765}-\frac{4A(h/2)}{651}+\frac{32A(h/4)}{279}
-\frac{512A(h/8)}{651}+\frac{16384A(h/16)}{9765}.
\end{equation}
For a term $c_ph^p$ the residual coefficient is multiplied by the exact moment
$M_p=\sum_{j=0}^{m}w_j2^{-jp}$. Direct rational evaluation gives
\begin{equation}\label{eq:moments}
M^{\{2,3\}}_4=\frac1{56},\qquad M^{\{2,3,4\}}_5=-\frac1{960},\qquad M^{\{2,3,4,5\}}_6=\frac1{31744}.
\end{equation}
Thus, when the indicated next coefficient is nonzero and a sufficiently long expansion holds,
the three formulas have leading residuals $c_4h^4/56$, $-c_5h^5/960$, and $c_6h^6/31744$,
respectively, and leading orders 4, 5, and 6. The sums of absolute weights are approximately
2.143, 2.429, and 2.585, respectively: these give modest worst-case amplification factors for
independent bounded input errors, although severe cancellation can still demand high working
precision. They require more refinement levels than the even-only operators to reach a
comparable order, and do not guarantee improved wall-clock efficiency. Numerical conditioning
and computational cost should be measured rather than presumed.

\paragraph{Numerical check on the perturbed family.}
In the test of Section~\ref{sec:numtest}, the operator \eqref{eq:P23} applied to the
scale-dependent family with $\varepsilon=0.1$ gives observed orders $4.008$ ($N=32\to64$),
$4.006$ ($N=64\to128$) and $4.003$ ($N=128\to256$), against third order for every even-power operator on the same data.
For $\varepsilon=0.2$ the corresponding values are $3.171$, $3.758$ and $3.903$, still rising
toward four within this window; the larger amplitude has a longer pre-asymptotic range.

%------------------------------------------------------------
\section{Spherical cap area: a geodesic-polygon family}\label{sec:cap}
This family is not a flat-triangle rule and is not covered by Theorem~\ref{thm:classical}.
For a cap with polar angle $\alpha$ on the unit sphere, $A_{\mathrm{cap}}=2\pi(1-\cos\alpha)$.
Join the north pole to adjacent vertices of the regular geodesic boundary polygon and use the
exact spherical area of each triangle~\cite{todhunter}. With $c=\cos\alpha$, $s=\sin\alpha$,
and $h=2\pi/N$, the oriented spherical triangle area for sufficiently small $h>0$ is
\begin{equation}\label{eq:E}
E_\alpha(h)=2\arctan\!\left(\frac{s^{2}\sin h}{1+2c+c^{2}+s^{2}\cos h}\right).
\end{equation}
(This follows from $\tan(E/2)=t\sin h/(1+t\cos h)$ with $t=\tan^2(\alpha/2)=(1-c)/(1+c)$,
after multiplying numerator and denominator by $(1+c)^2/(1+c)$ and using $1-c^2=s^2$.)
For $0<\alpha<\pi$ the denominator is positive near zero and even in $h$, while the numerator
is odd. Therefore $E_\alpha$ is odd analytic near zero:
\[
E_\alpha(h)=a_1h+a_3h^3+a_5h^5+\cdots,\qquad a_1=1-\cos\alpha .
\]
Since $A_N=NE_\alpha(2\pi/N)$, the polygon area has an even-power error expansion, and its
constant term $2\pi a_1=2\pi(1-\cos\alpha)$ equals the cap area exactly. This is the
one-dimensional mechanism familiar from polygonal approximation of the circle; it is a proof
for this cap family only.

\begin{table}[ht]
\centering
\begin{tabular}{lrrrr}
\toprule
$\alpha$ & Raw & Two-scale & Three-scale & Four-scale\\
\midrule
$15^\circ$ & 1.99989583 & 3.99999065 & 6.00093925 & 7.99993479\\
$30^\circ$ & 1.99996740 & 4.00041698 & 6.00000053 & 7.95025527\\
$60^\circ$ & 2.00016296 & 4.00014936 & 6.00003853 & 7.99880645\\
$75^\circ$ & 2.00023455 & 4.00028432 & 6.00028420 & 8.00026989\\
\bottomrule
\end{tabular}
\caption{Original cap experiment: observed orders formed by comparing $N=128$ with $N=256$
(the row label is the cap polar angle $\alpha$). The $30^\circ$ four-scale value $7.95025527$
is reproduced at 80-digit precision and represents a pre-asymptotic deviation, not a precision
floor.}\label{tab:cap}
\end{table}

%------------------------------------------------------------
\section{Latitude--longitude cell area: an independent direct proof}\label{sec:latlon}
For the unit-sphere cell $\lambda_0\le\lambda\le\lambda_1$ and $\phi_0\le\phi\le\phi_1$, the
exact area is
\[
A_{\mathrm{cell}}=(\lambda_1-\lambda_0)(\sin\phi_1-\sin\phi_0).
\]
The approximation uses a uniform latitude--longitude parameter grid and two Euclidean chordal
triangles per parameter cell. By Corollary~\ref{cor:latlon} this family is even-asymptotic to
every order. Proposition~\ref{prop:latlon-even} below, written before the author knew of
\cite{verlindencools,lyness94}, proves the same statement directly; it is retained as an
independent elementary argument that uses only cell parity and the midpoint Euler--Maclaurin
formula on a rectangle.

\begin{table}[ht]
\centering
\begin{tabular}{lrrrr}
\toprule
Cell & Raw & Two-scale & Three-scale & Four-scale\\
\midrule
Equatorial $30^\circ\times30^\circ$ & 1.99996836 & 3.99999637 & 5.99994223 & 7.99911861\\
Midlatitude $30^\circ\times30^\circ$ & 1.99998486 & 3.99996146 & 5.99998582 & 8.00085300\\
Wide equatorial $60^\circ\times30^\circ$ & 1.99992674 & 3.99992210 & 6.00000864 & 8.00036842\\
Northern $15^\circ\times60^\circ$ & 1.99994200 & 3.99996772 & 5.99996443 & 8.00007636\\
\bottomrule
\end{tabular}
\caption{Original latitude--longitude experiment: observed orders at $N=64$.}\label{tab:latlon}
\end{table}

\begin{proposition}[Even-power expansion for uniform latitude--longitude chordal cells]
\label{prop:latlon-even}
Fix $-\pi/2<\phi_0<\phi_1<\pi/2$ and $\lambda_0<\lambda_1$.
Subdivide the corresponding unit-sphere parameter rectangle into a uniform
$N\times N$ grid, with every cell triangulated along its $(p_{00},p_{11})$
diagonal, and let $A_N$ be the sum of the unsigned Euclidean areas of its two
chordal triangles. With $h=1/N$, for every positive integer $r$ there are
coefficients $c_{2},\ldots,c_{2r}$, independent of $N$, such that
\begin{equation}\label{eq:latlon-even}
 A_N=(\lambda_1-\lambda_0)(\sin\phi_1-\sin\phi_0)
       +\sum_{k=1}^{r}c_{2k}h^{2k}+O(h^{2r+2}).
\end{equation}
The assertion concerns this fixed rectangular domain and uniform, consistently
triangulated refinement; it does not extend to arbitrary perturbed grids.
\end{proposition}
\begin{proof}[Direct proof]
Put $\Delta\phi=\phi_1-\phi_0$ and $\Delta\lambda=\lambda_1-\lambda_0$.
A parameter cell has centre $(x,y)$ and corners
$(x\pm\Delta\phi\,h/2,y\pm\Delta\lambda\,h/2)$.
For the fixed $p_{00}$--$p_{11}$ diagonal, let $g(x,y,h)$ be the sum of the
unsigned Euclidean areas of the chordal triangles
$(p_{00},p_{10},p_{11})$ and $(p_{00},p_{11},p_{01})$.
Changing the sign of $h$ permutes the opposite corners; taken together, the
two unsigned triangle areas are invariant under this relabelling. Their sum therefore satisfies
$g(x,y,-h)=g(x,y,h)$ for real $h$ sufficiently close to zero.
Each triangle has area one half of the norm of a cross product of edge
vectors analytic in $h$. After factoring $h^2$, each squared norm has a
positive leading coefficient proportional to
$(\Delta\phi\,\Delta\lambda\cos x)^2$; hence its positive square root
is real analytic for small signed $h$, uniformly on the closed latitude band. Thus $g/h^2$
extends to an even real-analytic function of $h$, uniformly with respect to $(x,y)$ on the
fixed rectangle. In particular, for any fixed $r$,
\[
 g(x,y,h)=\sum_{k=1}^{r+1}\rho_{2k}(x,y)h^{2k}
                +O(h^{2r+4}),\qquad
 \rho_2(x,y)=\Delta\phi\,\Delta\lambda\cos x,
\]
with a uniform remainder. Rotational invariance makes the coefficient
functions independent of longitude $y$, although this is not required for
the midpoint argument.
Summing the remainder over $N^2=h^{-2}$ cells gives
$O(h^{2r+2})$. For each coefficient, the cell-centre sum is a composite
midpoint quadrature on a fixed rectangle, and the Euler--Maclaurin formula
for smooth coefficient functions has an expansion in even powers of $h$.
Collecting terms yields \eqref{eq:latlon-even}. The leading midpoint integral
is
\[
 \int_{\phi_0}^{\phi_1}\int_{\lambda_0}^{\lambda_1}
            \cos\phi\,d\lambda\,d\phi
 = (\lambda_1-\lambda_0)(\sin\phi_1-\sin\phi_0).
\]
\end{proof}

\begin{remark}[Scope of the parity argument]
The midpoint nodes and the fixed diagonal are part of the definition of the
refinement family. The assertion is an asymptotic expansion to each finite
order, not a claim that an infinite series converges. In particular it does
not cover the scale-dependent interior-node perturbation in Section~\ref{sec:perturb}. The
pairing of the two triangles of a cell is what makes the local term even here; for a single
triangle family the local term is not even, which is the difficulty identified
in~\cite{georgtausch} and resolved in~\cite{verlindencools,lyness94}.
\end{remark}

%------------------------------------------------------------
\section{Illustration of the classical theorem I: chordal spherical triangle}\label{sec:triangle}
Take $A=(0,0,1)$, $B=(1,0,0)$, and $C=(\cos\beta,\sin\beta,0)$. This is a birectangular
spherical triangle: the sides $AB$ and $AC$ are meridians and $BC$ lies on the equator, so the
angles at $B$ and $C$ are right angles and the angle at the pole $A$ is $\beta$. Its spherical
excess is $\beta+\pi/2+\pi/2-\pi=\beta$, so $\Aex=\beta$. Map a normalized barycentric
triangulation to the sphere and sum the Euclidean areas of its chordal triangles. Here
$\det[A,B,C]=\sin\beta\neq0$ for $0<\beta<\pi$, so Corollary~\ref{cor:triangle} applies and the
orders in Table~\ref{tab:triangle} illustrate a proved expansion. As a check of
\eqref{eq:jac}, $\iint_\sigma\sin\beta\,\|P\|^{-3}\,du\,dv=\beta$ was confirmed numerically to
50 digits for the four tabulated angles.

\begin{table}[ht]
\centering
\begin{tabular}{lrrrr}
\toprule
$\beta$ & Raw & Two-scale & Three-scale & Four-scale\\
\midrule
$15^\circ$ & 1.99942072 & 3.99915618 & 5.99886117 & 7.99716985\\
$30^\circ$ & 1.99938698 & 3.99908575 & 5.99880136 & 7.99761444\\
$60^\circ$ & 1.99923969 & 3.99881849 & 5.99858915 & 7.99835976\\
$90^\circ$ & 1.99885313 & 3.99827889 & 5.99785193 & 7.99808456\\
\bottomrule
\end{tabular}
\caption{Original chordal triangle experiment: observed orders at $N=64$.}\label{tab:triangle}
\end{table}

%------------------------------------------------------------
\section{Illustration of the classical theorem II: irregular spherical polygons}\label{sec:polygon}
The original study considered fixed irregular convex spherical polygons triangulated from a
root vertex. Each fan triangle was refined using a normalized barycentric chordal patch. One
six-vertex example had angular vertices $(4^\circ,-8^\circ)$, $(18^\circ,5^\circ)$,
$(15^\circ,28^\circ)$, $(2^\circ,42^\circ)$, $(-14^\circ,30^\circ)$, and $(-10^\circ,4^\circ)$;
the reported reference area by spherical excess was
$0.325062777898240815426301649532$. By Corollary~\ref{cor:fan}, an irregular but fixed
boundary cannot prevent the ladder, provided each fan triangle is non-degenerate; what the
boundary geometry controls is whether the fan tiles the intended region, that is, whether the
constant term is the target area.

\begin{table}[ht]
\centering
\begin{tabular}{rrrrr}
\toprule
$N$ & Raw & Two-scale & Three-scale & Four-scale\\
\midrule
16 & 1.99929248 & 3.99269841 & 5.99184034 & 8.01301620\\
32 & 1.99982228 & 3.99817005 & 5.99798552 & 8.00310155\\
64 & 1.99995552 & 3.99954224 & 5.99949790 & 8.00076594\\
\bottomrule
\end{tabular}
\caption{Original fixed six-vertex polygon experiment.}\label{tab:poly6}
\end{table}

\begin{table}[ht]
\centering
\begin{tabular}{rrrrr}
\toprule
Vertices & Raw & Two-scale & Three-scale & Four-scale\\
\midrule
5 & 1.99993053 & 3.99981653 & 5.99991832 & 7.99955672\\
6 & 1.99993586 & 3.99979158 & 5.99991390 & 8.00029966\\
8 & 2.00001032 & 4.00140962 & 5.99925770 & 7.99912295\\
\bottomrule
\end{tabular}
\caption{Original aggregate final-level orders over deterministic polygon families and
seeds.}\label{tab:polybatch}
\end{table}

\paragraph{Boundary stress sweep.}
The original polygon tests separated fixed irregular boundaries from changes in the mesh rule.
In a further boundary stress sweep, five-vertex cases at radii $15^\circ$ and $30^\circ$ with
jitter 0.2 and 0.5, and eight-vertex cases at jitter 0.2, retained near-$2,4,6,8$ behavior.
Eight-vertex jitter-0.5 cases did not converge to the same reference area, with raw errors
nearly constant and observed orders near zero. Corollary~\ref{cor:fan} supports the earlier
reading of these cases: each fan triangle still carries an even-power expansion, so a nearly
constant error indicates that the fan's constant term differs from the reference area
(orientation, nonconvexity, overlap, or a different limiting region), not a failure of the
expansion. This diagnosis has not been confirmed case by case. For all test families, the
domain, orientation, target area, and approximation rule must be specified unambiguously.

%------------------------------------------------------------
\section{Limitations and validation priorities}
The even-power structure is now established for every fixed-map family in this manuscript: by
the cap parity argument, and by Theorem~\ref{thm:classical} for the chordal triangle, fan
polygon and latitude--longitude families. The tabulated observed orders remain numerical
illustrations attributed to the original manuscript. The original tests use convex or
otherwise well-behaved spherical domains with available reference areas. Adaptive,
scale-dependent and highly irregular meshes, and fans that do not tile the intended domain,
are outside the established scope. The hypothesis of Proposition~\ref{prop:cubic} is proved (Theorem~\ref{thm:joint}) only for
uniform square-cell grids with a fixed diagonal, deformed by a fixed smooth field that
vanishes on the boundary. Independent node-by-node jitter, which is not the image of a smooth
field, and adaptive or topology-changing triangulations are not covered, and the nonvanishing
of $C_2'(0)$ is certified for one field on one cell only. The proofs of
Theorem~\ref{thm:joint} and Lemma~\ref{lem:q2} and the certificate of
Proposition~\ref{prop:certified} were produced with AI assistance; they have been checked by
exact computer algebra and by reruns on the author's machine, but have not had an independent
human proof reading, and no claim of historical novelty is established by the limited
literature search reported in Section~\ref{sec:scope}. The numerical test of
Section~\ref{sec:numtest} has not been repeated on other fields or geometries. The numerical anomalies of the original perturbed-grid
Table~\ref{tab:perturb} are \emph{not} reproduced by this revision, because the original field
and seeds are unavailable; the cubic mechanism explains them in kind only. Supplementary
AI-assisted diagnostic fit tables requiring missing node and precision metadata are omitted
from this release. Section~\ref{sec:numtest} carries out items (i), (ii), (iv) and (v) of the following programme
for a single field, and item (iii) in part; item (vi) is open. To strengthen the numerical
result further, one should (i) test fixed-map and $h$-dependent-map meshes with the same node field, (ii) retain
boundary constraints and triangle orientation checks, (iii) sweep several $N$ and deformation
amplitudes, (iv) fit the $h^3$ coefficient directly, (v) verify precision stability, and
(vi) compare both error-versus-cost and observed-order-versus-cost across even-only and
mixed-power operators.

%------------------------------------------------------------
\section{Relation to the general N-Series framework}
Spherical area supplies another setting in which geometrically defined approximants exhibit
extrapolatable error modes. For flat-triangle approximations of smooth surfaces this was
established three decades ago in the boundary-element literature; for those fixed-map families the present manuscript adds spherical specializations and
examples, not mathematical novelty in the extrapolation or in the expansion. Its own
contribution lies at the boundary of that theory: the cubic obstruction for scale-dependent
families and the mixed-exponent operators that remove it. The episode also illustrates a general point for the N-Series programme: when a
geometric approximant shows an even-power ladder under regular refinement, the
Euler--Maclaurin literature for the relevant element shape should be checked before the
ladder is described as conjectural. Williamson~\cite{williamson} surveys the regular
spherical grids on which such area and quadrature questions arise in practice;
Coxeter~\cite{coxeter} is a general reference for the spherical geometry used.

%------------------------------------------------------------
\section{Conclusion}
The $2,4,6,8$ ladder for spherical area is a property of fixed smooth refinement families,
and for chordal triangulations it is classical. Its boundary is the subject of this
manuscript. A displacement proportional to the grid spacing takes a family outside the
classical theorem. For uniform square-cell grids deformed by a fixed smooth field we prove an
even-power expansion that is uniform in the displacement, from which the scale-dependent
family acquires a cubic term
$\varepsilon C_2'(0)h^3$ that even-power extrapolation cannot remove; the two-, three- and
four-scale operators leave it with exact factors $-1/6$, $-1/90$ and $-1/810$. For the test
field the coefficient $C_2'(0)$ is given by an explicit integral and certified negative by
interval arithmetic. A
deterministic test confirms the amplitude law, the factor ratios $1/15$ and $1/9$, the stall
of the extrapolated orders at three, and the recovery of fourth order by the mixed-exponent
operator for $\{2,3\}$. Mixed-exponent operators with exact weights and moments are available
whenever the exponents are identified. For fixed maps, the even-power expansion is proved
here for spherical caps bounded by regular geodesic polygons, and follows from the theorem of
Verlinden--Cools and Lyness, proving a conjecture of Georg, for chordal spherical triangles,
fan-refined polygons and uniform latitude--longitude cells. The open items are an independent human reading of the new proofs and of the certificate, a
fuller search of the literature on perturbed meshes, an extension from smooth fixed-field
deformations to independent node jitter, tests on further fields and geometries, and a
reconstruction of the original perturbed-grid experiment, whose anomalies are explained in
kind but not reproduced.

%------------------------------------------------------------
\appendix
\section{Reproducibility notes for the original experiments}
The original June 2026 manuscript identifies the following standalone Python scripts:
\begin{itemize}
\item \texttt{spherical\_area\_nseries\_cap\_test.py}
\item \texttt{spherical\_triangle\_nseries\_area\_test.py}
\item \texttt{spherical\_latlon\_cell\_nseries\_area\_test.py}
\item \texttt{spherical\_irregular\_polygon\_nseries\_test.py}
\item \texttt{spherical\_irregular\_polygon\_batch\_nseries\_test.py}
\item \texttt{spherical\_irregular\_polygon\_stress\_sweep\_nseries\_test.py}
\item \texttt{spherical\_perturbed\_latlon\_nseries\_test.py}
\end{itemize}
For errors $e_N=|A_N-\Aex|$, the observed order is $p_N=\log(e_N/e_{2N})/\log2$; the
corresponding definition applies to the extrapolated errors. Some original experiment scripts
were not included in the source/PDF review package, and their archived execution status is not
established here. Numerical tables carried over from earlier versions are attributed to the
original manuscript; this statement is not a reproducibility certification.

\sloppy The test of Section~\ref{sec:numtest} is \texttt{scripts/cubic\_amplitude\_test\_mpfr.py}
(requires \texttt{gmpy2}); it is deterministic, platform-independent, performs the orientation
checks, and prints every number quoted there. Its output of November~10, 2026 is archived in
\texttt{output/}. The earlier \texttt{scripts/cubic\_amplitude\_test.py} uses the platform's
\texttt{long double} (x86 80-bit in the archived run) and serves as a second implementation. The statements of Section~\ref{sec:jointthm} are supported by
\texttt{scripts/q2\_lemma\_symbolic\_audit.py} (exact polynomial audit of
Lemma~\ref{lem:q2}; SymPy), \texttt{scripts/c2prime\_enclosure.py} (certified enclosure of
Proposition~\ref{prop:certified}; SymPy and python-flint),
\texttt{scripts/c2prime\_crosscheck\_numpy.py} (the same integrals by Gauss--Legendre
quadrature; no interval library needed) and \texttt{scripts/check\_joint\_expansion.py}
(formula \eqref{eq:C2formula} against grid sums). Their outputs of November~10, 2026 are
archived in \texttt{output/}. The other checks
added in this draft are in \texttt{scripts/check\_v2\_4\_draft.py}: the Jacobian
integral of Lemma~\ref{lem:radial} for the four tabulated triangles; the octant errors of
Table~\ref{tab:gt}, computed with the archived routine \texttt{chordal\_triangle}; and exact
rational evaluation of the weights \eqref{eq:R2}--\eqref{eq:R4} and
\eqref{eq:P23}--\eqref{eq:P2345}, the cubic residual factors and the moments
\eqref{eq:moments}. The correspondence between the archived routine and the regular
subdivision of Section~\ref{sec:classical} was established by reading the routine's index
ranges.

\section{Revision history and provenance}
\textbf{Version 1 (June 3, 2026):} five spherical area tests, tables of observed orders, cap
parity argument, classical even-power scale operators, and limitations.

\textbf{Version 2 (October 24, 2026):} precise expansion-based regularity definition;
conditional $h^3$ perturbation proposition; AI-assisted exploratory diagnostics, distinctly
labeled; generalized mixed-power operators; stronger limits on what numerical observations
establish. The original numerical values are preserved rather than silently replaced.

\textbf{Version 2.1 (October 24, 2026):} corrected the regularity needed for the $O(h^5)$
remainder; clarified provenance as exploratory AI-assisted review; added first surviving exact
moments, error amplification factors, and a falsifiable signed-amplitude test of
cubic-coefficient linearity.

\textbf{Review supplement (October 24, 2026):} added specified deformation fields,
amplitude-fit sensitivity and provenance distinctions, plus an AI-assistance disclosure. The
original perturbation experiment remains unverified.

\textbf{Version 2.2 (October 25, 2026):} stated the exact cubic residual factors $-1/6$,
$-1/90$, $-1/810$ with a product-of-stages form of the operators; rewrote the cap formula
numerator as $s^2\sin h$ with a one-line derivation; described the test triangle as
birectangular; described the cubic mechanism as consistent with, not established by, the
perturbed-grid table; omitted supplementary fit data with incomplete provenance. No numerical
table values were changed.

\textbf{Version 2.3 release candidate (October 29, 2026):} added
Proposition~\ref{prop:latlon-even} with its direct proof. Updated the abstract, scope and
conclusion without changing the original numerical tables.

\textbf{Version 2.4 draft (November 10, 2026):} reorganized to lead with the manuscript's own
results (the cubic obstruction, its numerical confirmation, the mixed-exponent operators and
the cap family) and to present as background the literature on
flat-triangle approximation of surface integrals
\cite{georg91,georgtausch,verlindencools,lyness94}, which earlier versions did not cite.
Added Theorem~\ref{thm:classical} (cited), Lemma~\ref{lem:radial},
Corollaries~\ref{cor:triangle}--\ref{cor:latlon}, Remarks~\ref{rem:gt} and~\ref{rem:supplies},
Table~\ref{tab:gt}, and the new numerical test of Section~\ref{sec:numtest} with
Tables~\ref{tab:b3} and~\ref{tab:ratios}. A later pass on the same day folded in a proof supplement:
Theorem~\ref{thm:joint}, Lemma~\ref{lem:q2}, Corollary~\ref{cor:cubicfamily} and
Proposition~\ref{prop:certified}, with a related-work paragraph and the certified value of
$C_2'(0)$ in place of a coarse finite-difference estimate ($-6.806\times10^{-4}$) used in the
first pass. No previously stated theorem, proposition, lemma or corollary was altered in
either pass. The chordal triangle and fan-polygon families are reclassified from
numerical evidence to instances of a classical theorem; Proposition~\ref{prop:latlon-even} is
reclassified from a new result to an independent proof of a known one. Earlier versions made
no claim of a general theorem, so no previously stated result is withdrawn; what changes is
attribution and status. Equation and table numbering no longer match earlier versions. No
carried-over numerical table values were changed.

\paragraph{AI assistance in this revision.}
The relevant literature was identified, its hypotheses were compared with the archived
implementation, the numerical test of Section~\ref{sec:numtest} was designed, coded in two implementations and run, and this
restructured draft was prepared with the assistance of Claude (Anthropic) at the author's
direction. The proof of Theorem~\ref{thm:joint} was first drafted as a separate AI-assisted technical
supplement at the author's direction and then audited and tightened in this session; the
proof of Lemma~\ref{lem:q2}, its symbolic audit and the interval enclosure were produced in
this session. A separately written AI-assisted direct-grid computation corroborated the sign
and magnitude of $C_2'(0)$. The author reran \texttt{c2prime\_crosscheck\_numpy.py} and
\texttt{c2prime\_enclosure.py} on his own Windows machine on November~10, 2026 and obtained the
same values, including the certified sign. Both implementations of the numerical test came from the same AI-assisted session, so their agreement shows
precision stability, not independence. The scripts are archived and deterministic so that the
author and readers can rerun them; an independent rerun is pending. The 1994 papers
\cite{georgtausch,verlindencools,lyness94} were read in full for this purpose;
\cite{georg91} was not available and is cited on the authority of the three papers that state
and prove its Conjecture~5.1. This is an AI-assisted revision, not an independent human peer
review, and the author's own review of this draft is pending.

%------------------------------------------------------------
\begin{thebibliography}{99}
\bibitem{austintrefethen} A. P. Austin and L. N. Trefethen, ``Trigonometric interpolation and
quadrature in perturbed points,'' \emph{SIAM Journal on Numerical Analysis}, vol.~55, no.~5,
2017, pp.~2113--2122.
\bibitem{blumlinrannacher} H. Blum, Q. Lin, and R. Rannacher, ``Asymptotic error expansion and
Richardson extrapolation for linear finite elements,'' \emph{Numerische Mathematik}, vol.~49,
1986, pp.~11--37.
\bibitem{brs} F. L. Bauer, H. Rutishauser, and E. Stiefel, ``New aspects in numerical
quadrature,'' in \emph{Proceedings of Symposia in Applied Mathematics}, vol.~15, American
Mathematical Society, 1963, pp.~199--218.
\bibitem{bulirschstoer} R. Bulirsch and J. Stoer, ``Fehlerabsch\"atzungen und Extrapolation
mit rationalen Funktionen bei Verfahren vom Richardson-Typus,'' \emph{Numerische Mathematik},
vol.~6, 1964, pp.~413--427.
\bibitem{coxeter} H. S. M. Coxeter, \emph{Introduction to Geometry}, Wiley, 1969.
\bibitem{davisrabinowitz} P. J. Davis and P. Rabinowitz, \emph{Methods of Numerical
Integration}, 2nd ed., Academic Press, 1984.
\bibitem{diskinthomas} B. Diskin and J. L. Thomas, ``Effects of mesh irregularities on accuracy
of finite-volume discretization schemes,'' AIAA paper, 2012; NASA Technical Reports Server
document 20120001451.
\bibitem{georg91} K. Georg, ``Approximation of integrals for boundary element methods,''
\emph{SIAM Journal on Scientific and Statistical Computing}, vol.~12, no.~2, 1991,
pp.~443--453.
\bibitem{georgtausch} K. Georg and J. Tausch, ``Some error estimates for the numerical
approximation of surface integrals,'' \emph{Mathematics of Computation}, vol.~62, no.~206,
1994, pp.~755--763.
\bibitem{johansson} F. Johansson, ``Arb: efficient arbitrary-precision midpoint-radius interval
arithmetic,'' \emph{IEEE Transactions on Computers}, vol.~66, no.~8, 2017, pp.~1281--1292.
\bibitem{lyness78} J. N. Lyness, ``Quadrature over a simplex. II. A representation for the
error functional,'' \emph{SIAM Journal on Numerical Analysis}, vol.~15, no.~5, 1978,
pp.~870--887.
\bibitem{lyness94} J. N. Lyness, ``Quadrature over curved surfaces by extrapolation,''
\emph{Mathematics of Computation}, vol.~63, no.~208, 1994, pp.~727--740.
\bibitem{marchukshaidurov} G. I. Marchuk and V. V. Shaidurov, \emph{Difference Methods and
Their Extrapolations}, Springer, 1983.
\bibitem{richardson} L. F. Richardson, ``The approximate arithmetical solution by finite
differences of physical problems,'' \emph{Philosophical Transactions of the Royal Society A},
vol.~210, 1911, pp.~307--357.
\bibitem{sidi} A. Sidi, \emph{Practical Extrapolation Methods: Theory and Applications},
Cambridge University Press, 2003.
\bibitem{todhunter} I. Todhunter, \emph{Spherical Trigonometry}, Macmillan, 1886.
\bibitem{verlindencools} P. Verlinden and R. Cools, ``Proof of a conjectured asymptotic
expansion for the approximation of surface integrals,'' \emph{Mathematics of Computation},
vol.~63, no.~208, 1994, pp.~717--725.
\bibitem{williamson} D. L. Williamson, ``The Evolution of Dynamical Cores for Global
Atmospheric Models,'' \emph{Journal of the Meteorological Society of Japan}, vol.~85B, 2007,
pp.~241--269.
\end{thebibliography}

\end{document}
